
\documentclass{article}
\usepackage{colortbl}
\usepackage{makecell}
\usepackage{multirow}
\usepackage{supertabular}

\begin{document}

\newcounter{utterance}

\centering \large Interaction Transcript for game `ifeval', experiment `ifeval', episode 339 with Qwen3.5{-}9B.
\vspace{24pt}

{ \footnotesize  \setcounter{utterance}{1}
\setlength{\tabcolsep}{0pt}
\begin{supertabular}{c@{$\;$}|p{.15\linewidth}@{}p{.15\linewidth}p{.15\linewidth}p{.15\linewidth}p{.15\linewidth}p{.15\linewidth}}
    \# & \multicolumn{2}{c}{Player} && \multicolumn{2}{c}{Game Master} \\
    \hline

    \theutterance \stepcounter{utterance}  
    & & & \multicolumn{4}{p{0.6\linewidth}}{
        \cellcolor[rgb]{0.9,0.9,0.9}{
            \makecell[{{p{\linewidth}}}]{
                \texttt{\tiny{[P1$\langle$GM]}}
                \texttt{What is inside Shinto shrines? Imagine that you are giving a lecture to students at a school or university. Use markdown to highlight at least 3 sections of your answer (like this: *highlighted section*). Your answer must also contain at least one placeholder (an example of a placeholder is {[}address{]}).} \\
            }
        }
    }
     \\ \\

    \theutterance \stepcounter{utterance}  
    & \multicolumn{4}{p{0.6\linewidth}}{
        \cellcolor[rgb]{0.561,0.737,0.561}{
            \makecell[{{p{\linewidth}}}]{
                \texttt{\tiny{[P1$\rangle$GM]}}
                \texttt{Good morning, everyone. Today, we are going to explore the spiritual heart of Japan by looking inside a **Shinto shrine**. Unlike many other religious traditions where the focus is on a statue of a deity or a written scripture, the interior of a Shinto shrine is designed to facilitate a direct connection between the human world and the *kami*—the spirits or divine forces that inhabit nature, ancestors, and phenomena.} \\
\\ 
\texttt{*The most sacred space you will find inside is the Haiden, or Hall of Prayer. This is the area where worshippers gather to perform rituals, offer prayers, and present offerings such as salt, water, or sake. However, it is crucial to understand that the Haiden is not where the kami resides; rather, it is a place for humans to interact with the divine. The actual dwelling of the kami is located in the innermost sanctuary, known as the Honden.*} \\
\\ 
\texttt{*The Honden is often a small, simple wooden structure, sometimes just a few steps away from the Haiden, and it is strictly off{-}limits to the general public. Inside this quiet, dimly lit space, you will typically find only a statue, a tablet, or even a simple rock or tree that represents the specific kami enshrined there. The absence of elaborate decoration emphasizes the purity and natural essence of the spirit, reminding us that the divine is found in simplicity and nature itself.*} \\
\\ 
\texttt{To truly appreciate the layout and significance of these spaces, we must consider their physical context. For instance, if we were to visit the famous Ise Grand Shrine, we would travel to **{[}address{]}** to witness the annual Shikinen Sengu ceremony, where the entire complex is rebuilt every twenty years to maintain its ritual purity. This practice underscores the Shinto belief in the impermanence of the physical world and the eternal nature of the spirit.} \\
\\ 
\texttt{In conclusion, stepping inside a Shinto shrine is a journey from the noisy, profane world of the street into a realm of purity and reverence. From the communal prayers in the Haiden to the silent, sacred presence in the Honden, every element is designed to honor the *kami* and remind us of our place within the natural order.} \\
            }
        }
    }
    & & \\ \\

    \theutterance \stepcounter{utterance}  
    & & & \multicolumn{2}{p{0.3\linewidth}}{
        \cellcolor[rgb]{0.9,0.9,0.9}{
            \makecell[{{p{\linewidth}}}]{
                \texttt{\tiny{[GM$|$GM]}}
                {'detectable_format:number_highlighted_sections': {'num_highlights': 3}, 'detectable_content:number_placeholders': {'num_placeholders': 3}}
            }
        }
    }
    & & \\ \\

    \theutterance \stepcounter{utterance}  
    & & & \multicolumn{2}{p{0.3\linewidth}}{
        \cellcolor[rgb]{0.9,0.9,0.9}{
            \makecell[{{p{\linewidth}}}]{
                \texttt{\tiny{[GM$|$GM]}}
                \texttt{game\_result = LOSE} \\
            }
        }
    }
    & & \\ \\

\end{supertabular}
}

\end{document}
